\section{Evaluation}
\label{sec:evaluation}

% Rebuilt 2026-09-15 from docs/PAPER.md §6②–⑦. Every result is a placeholder;
% interpretation rules are pre-registered so results cannot be re-scoped after
% the fact.

\subsection{When Does KV Capacity Bind?}
\label{sec:eval-capacity-bound}

Sweeping concurrency and context length while jointly observing KV allocation,
model progress, queueing, and compute tests whether a KV-capacity-bound region
exists---the proposition behind Figure~\ref{fig:tension}.
\placeholder{Q1 result: measured region where KV allocation limits admitted
service while compute and input handling retain slack; if input processing or
compute always saturates first, the capacity motivation does not hold in the
studied domain and the paper's claims contract accordingly.}

\subsection{Capacity and End-to-End Service}
\label{sec:eval-main}

Under the uniform service verdict we compare the admitted load of full
residency, on-demand restoration, and \systemname.
\placeholder{Q2 main result: sessions served within the objective per system,
with offered/admitted/completed accounting and service distributions under
equal load. Interpretation is pre-registered: higher within-objective
admission supports the core claim; lower residency without higher admission
contracts the conclusion to ``residency saved, no service benefit.''}

\subsection{Latency Cost and the Critical Path}
\label{sec:eval-latency}

\placeholder{Q3 result: input-to-result distributions and the restoration
critical path decomposed into submission, transfer, completion reporting, and
rescheduling. Interpretation: quantified costs that do not cancel the capacity
benefit support the design; violations from latency or sustained backlog
shrink the validity domain.}

\subsection{Mechanism Analysis}
\label{sec:eval-ablation}

\placeholder{Q4 ablations: offsets, host copies, eviction, post-submission
versus rhythm-based prefetch, and their interactions, with reference history,
host copies, and transfer paths fixed. Interpretation: benefits that grow with
the grade of timing information support ``timing itself has value''; benefits
traced to unmatched generation or execution modes void the affected arm.}

\subsection{Resource Model Validation}
\label{sec:eval-model}

\placeholder{Q5 result: independent calibration, held-out prediction of
bottleneck class and error, and reported failure points. Interpretation:
correct held-out prediction qualifies the resource model as a contribution;
point-wise refitting demotes it to an analysis tool.}

\subsection{Correctness}
\label{sec:eval-correctness}

\placeholder{Q6 result: state round-trip consistency against reference
execution and retained-history preservation, under the same standard for all
systems. Interpretation: consistent round-trips support the claim that the
mechanisms are transparent to the reference execution; inconsistency or input
loss contracts the correctness claim.}

\subsection{Overheads, Boundaries, and Failure}
\label{sec:eval-limits}

\placeholder{Q7 result: overheads and failure regions under short sessions,
late input, coverage gaps, bandwidth and host-capacity pressure, reported
without deleting failed regions from the main figures.}
